\documentclass[12pt,oneside]{book}

%==================================================================
% METADATOS
%==================================================================
\newcommand{\universidad}{Universidad de Buenos Aires (UBA)}
\newcommand{\facultad}{Facultad de Derecho}
\newcommand{\materia}{Contrato Inmobiliario}
\newcommand{\titulo}{Donación con pacto de reversión,\\ fideicomiso inmobiliario y tokenización}
\newcommand{\autor}{Isaac Edgar Camacho Ocampo}
\newcommand{\anio}{2026}

%==================================================================
% PAQUETES
%==================================================================
\usepackage[utf8]{inputenc}
\usepackage[T1]{fontenc}
\usepackage[spanish,es-nodecimaldot]{babel}

\usepackage[
    top=2.5cm,
    bottom=2.5cm,
    left=3cm,
    right=2.5cm,
    headheight=15pt
]{geometry}

\usepackage{amsmath}
\usepackage{amssymb}
\usepackage{mathtools}

\usepackage{graphicx}
\usepackage{float}
\usepackage{caption}
\usepackage{subcaption}

\usepackage{booktabs}
\usepackage{array}
\usepackage{longtable}
\usepackage{tabularx}

\usepackage{enumitem}

\usepackage{xcolor}
\usepackage[most]{tcolorbox}

\usepackage{fancyhdr}
\usepackage{ragged2e}

\usepackage[
    colorlinks=true,
    linkcolor=blue!50!black,
    citecolor=green!40!black,
    urlcolor=blue!60!black,
    pdftitle={Apuntes de Contrato Inmobiliario},
    pdfauthor={Isaac Edgar Camacho Ocampo},
    pdfsubject={Apuntes universitarios},
    bookmarks=true
]{hyperref}

%==================================================================
% CONFIGURACIÓN
%==================================================================
\graphicspath{
    {./img/}
    {./graficos/}
    {./figuras/}
}

\setcounter{tocdepth}{2}

\setlength{\parindent}{0pt}
\setlength{\parskip}{0.5em}

\setlist[itemize]{
    topsep=0.3em,
    itemsep=0.2em
}

\setlist[enumerate]{
    topsep=0.3em,
    itemsep=0.2em
}

\clubpenalty=10000
\widowpenalty=10000

\pagestyle{fancy}
\fancyhf{}
\fancyhead[L]{\nouppercase{\leftmark}}
\fancyhead[R]{\materia}
\fancyfoot[C]{\thepage}
\renewcommand{\headrulewidth}{0.4pt}
\renewcommand{\footrulewidth}{0pt}

\fancypagestyle{plain}{
    \fancyhf{}
    \fancyfoot[C]{\thepage}
    \renewcommand{\headrulewidth}{0pt}
}

%==================================================================
% COMANDOS
%==================================================================
\newcommand{\abs}[1]{\left|#1\right|}
\newcommand{\norm}[1]{\left\lVert#1\right\rVert}
\newcommand{\vect}[1]{\mathbf{#1}}
\newcommand{\deriv}[2]{\frac{d#1}{d#2}}
\newcommand{\pderiv}[2]{\frac{\partial#1}{\partial#2}}

%==================================================================
% ENTORNOS / CAJAS
%==================================================================
\newtcolorbox{definition}[1]{
    enhanced,
    breakable,
    title={Definición: #1},
    fonttitle=\bfseries,
    colback=blue!3,
    colframe=blue!50!black,
    boxrule=0.8pt,
    arc=2mm
}

\newtcolorbox{important}{
    enhanced,
    breakable,
    title={Importante},
    fonttitle=\bfseries,
    colback=yellow!8,
    colframe=orange!70!black,
    boxrule=0.8pt,
    arc=2mm
}

\newtcolorbox{example}[1]{
    enhanced,
    breakable,
    title={Ejemplo: #1},
    fonttitle=\bfseries,
    colback=green!4,
    colframe=green!40!black,
    boxrule=0.8pt,
    arc=2mm
}

\newtcolorbox{formula}[1]{
    enhanced,
    breakable,
    title={Fórmula / Regla: #1},
    fonttitle=\bfseries,
    colback=purple!4,
    colframe=purple!50!black,
    boxrule=0.8pt,
    arc=2mm
}

\newtcolorbox{exercise}[1]{
    enhanced,
    breakable,
    title={Ejercicio: #1},
    fonttitle=\bfseries,
    colback=gray!5,
    colframe=gray!60!black,
    boxrule=0.8pt,
    arc=2mm
}

\newtcolorbox{note}{
    enhanced,
    breakable,
    title={Nota},
    fonttitle=\bfseries,
    colback=white,
    colframe=gray!50,
    boxrule=0.6pt,
    arc=2mm
}

%==================================================================
% CONFIGURACIÓN FINAL
%==================================================================
\renewcommand{\arraystretch}{1.25}

%==================================================================
% DOCUMENTO
%==================================================================
\begin{document}

%------------------------------------------------------------------
% PORTADA
%------------------------------------------------------------------
\begin{titlepage}

\thispagestyle{empty}

\begin{center}

\vspace*{1cm}

{\large \textsc{\universidad}}

\vspace{0.2cm}

{\large \textsc{\facultad}}

\vspace{3cm}

{\Huge \textbf{\materia}}

\vspace{1cm}

{\LARGE \titulo}

\vspace{0.8cm}

\textit{Estudiar con método transforma la información en comprensión y la comprensión en criterio.}

\vspace{1.2cm}

\rule{0.70\textwidth}{0.5pt}

\vspace{0.5cm}

{\large Apuntes de estudio}

\vspace{0.5cm}

\rule{0.70\textwidth}{0.5pt}

\vfill

{\large \autor}

\vspace{0.3cm}

{\large \anio}

\end{center}

\vfill

\begin{flushleft}
\textit{Este es un modesto aporte a los estudiantes de ``Contrato Inmobiliario'' no pretende sustituir el material oficial de la materia.}
\end{flushleft}

\end{titlepage}

%------------------------------------------------------------------
% ÍNDICE
%------------------------------------------------------------------
\tableofcontents
\clearpage

%==================================================================
% CAPÍTULO 1
%==================================================================
\chapter[Criterio de evaluación y empresa familiar]{Apertura: criterio de evaluación y costado humano de la materia}

Los contenidos de este apunte se organizan en dos grandes bloques. El primero parte de una consulta sobre la \textbf{donación con usufructo y pacto de reversión}, el plazo de 10 años, la reacción del Registro de la Propiedad Inmueble, la planificación sucesoria y el régimen de observabilidad de las donaciones. El segundo desarrolla la \textbf{tokenización inmobiliaria}: el concepto de token, su relación con el fideicomiso, el \textit{white paper} y la tecnología \textit{blockchain}, el \textit{sandbox} normativo de la CNV, los proveedores de servicios de activos virtuales (PSAV), el ciclo de vida del token y el rol del escribano.

\section{Utilidad profesional de la materia}

Los contenidos se vinculan con la práctica profesional en tres sentidos que surgen del propio material de clase.

\textbf{Es lo que van a consultar los clientes.} Un cliente que desea donar un departamento a su hijo sin perder el control necesita que el profesional sepa redactar y defender un pacto de reversión y anticipar la posición del Registro. Un cliente que desea invertir en tokens necesita que se le indique si el proveedor está inscripto en la CNV (si no lo está, se trata de una estafa) y que el profesional sepa leer un \textit{white paper}.

\textbf{Enseña a decidir con criterio.} La evaluación no se centra en qué postura se elige, sino en si se la fundamenta en el Código y en la Constitución. Sostener una posición y justificarla es la habilidad central de un abogado o escribano.

\textbf{Tiene también un valor personal.} Comprender cómo se transmite un inmueble, cómo se planifica una herencia o qué es realmente un token protege de errores costosos, de esquemas Ponzi y de promesas de inversión que aparentan ser fáciles. Quien comprende jurídicamente la tokenización tiene, además, ventaja, dado que se trata de un tema de actualidad.

\section{Cómo se relacionan los temas: la analogía de la familia y el departamento}

El material propone un ejemplo cotidiano para vincular los bloques temáticos. Una familia desea que su hijo tenga un departamento. Los padres le donan parte del patrimonio e incluyen una cláusula para que, si el hijo fallece antes, el bien vuelva a ellos (donación y reversión). Un joven sin ahorros suficientes no puede comprar la unidad completa, por lo que accede a una plataforma donde adquiere \textbf{fracciones}, como quien junta figuritas para completar un álbum (tokenización). Para que no circulen «figuritas falsas», el Estado exige que el proveedor esté inscripto y regula el entorno (sandbox y PSAV). Al completar el álbum, el joven cambia sus fracciones por la unidad y el escribano redacta la escritura: cambia la causa, no las reglas (causa fuente del dominio). Lo que garantiza que cada fracción tenga dueño y fecha cierta es la tecnología \textit{blockchain}.

\section{La empresa familiar y el costado humano}

La materia incluye un bloque de \textbf{empresa familiar}, que se trabaja desde lo humano: allí la tecnología queda de lado y se aborda la dimensión de los vínculos. La razón de su importancia es que la empresa familiar suele ser el primer contacto que tendrá el profesional (escribano o abogado) para resolver cuestiones legales; las pymes, en general, son empresas familiares, y a cualquier profesional le toca, en algún momento, intervenir en una de ellas.

Existe abundante bibliografía y conflictos documentados sobre el tema. Se participa de la idea de que el profesional debe contar con herramientas de distintas disciplinas y conocer los vínculos y las relaciones humanas.

\section{Criterio de evaluación: aplicar el Código}

El rol de la cátedra es transmitir; el estudiante toma lo que puede. Se espera que el alumno se lleve conocimiento y también dudas, porque hay cuestiones que no pueden resolverse de manera definitiva, aunque existe la obligación de evaluar.

La orientación de estudio es seguir el \textbf{Código Civil y Comercial}: quien lo conoce siempre dispone de una respuesta. Se debe seguir el Código y el instituto de cada figura. Por ejemplo, si el tema es el fideicomiso, corresponde estudiar el contrato de fideicomiso completo en el Código Civil y Comercial, desde que empieza el artículo hasta que cierra el artículo. Ese conocimiento brinda el marco jurídico del instituto analizado y resulta suficiente; no es necesario leer una gran cantidad de textos.

Lo que se busca es que el estudiante \textbf{aplique} los institutos, más que devolver teoría. La teoría debe saberse, pero para saber aplicarla: en los casos planteados hay mucho conflicto en cada figura, y la existencia de una respuesta se garantiza con la lectura del Código. No se juzga la respuesta, sino la \textbf{calidad del conocimiento}. Quien adopta una postura opuesta a la de otra docente, pero la fundamenta jurídicamente, es valorado positivamente.

\begin{important}
\textbf{Criterio de evaluación:} no se evalúa la respuesta en sí, sino la calidad del conocimiento y el fundamento jurídico. Quien conoce el Código Civil y Comercial y aplica el instituto completo (del primer al último artículo) siempre tiene una respuesta, incluso si sostiene una postura distinta a la de la otra docente.
\end{important}

\section{Cuadro Cornell del capítulo}

\begin{longtable}{
    >{\raggedright\arraybackslash}p{0.27\textwidth}
    >{\raggedright\arraybackslash}p{0.65\textwidth}
}
\toprule
\textbf{Preguntas / palabras clave} & \textbf{Notas / respuestas} \\
\midrule
\endfirsthead
\toprule
\textbf{Preguntas / palabras clave} & \textbf{Notas / respuestas} \\
\midrule
\endhead
\bottomrule
\endfoot

\textbf{¿Cómo se evalúa la materia?} &
Se evalúa la calidad del conocimiento y el fundamento jurídico, no la respuesta en sí. Alcanza con conocer el Código Civil y Comercial completo y aplicar cada instituto de principio a fin; se puede sostener una postura distinta a la de la otra docente si se fundamenta. \\[0.5em]

\textbf{¿Por qué el profesional debe entender los vínculos humanos?} &
Porque la empresa familiar suele ser el primer contacto del profesional con estas cuestiones legales, y porque el profesional debe contar con herramientas de distintas disciplinas y conocer los vínculos y las relaciones humanas. \\[0.5em]

\textbf{¿Cuál es la utilidad profesional de los temas de la clase?} &
Permite asesorar a clientes que desean donar sin perder el control (pacto de reversión y posición del Registro) o invertir en tokens (verificar la inscripción del proveedor en la CNV y leer un \textit{white paper}), y entrena la habilidad de sostener y justificar una posición. \\

\midrule
\multicolumn{2}{p{0.92\textwidth}}{
\textbf{Resumen:} La materia se evalúa por la calidad del conocimiento y por el fundamento jurídico, tomando el Código Civil y Comercial como base completa de cada instituto. El objetivo es aplicar, no repetir teoría, y comprender también la dimensión humana de los casos, especialmente en la empresa familiar.
} \\
\end{longtable}

%==================================================================
% CAPÍTULO 2
%==================================================================
\chapter[Donación con usufructo y pacto de reversión]{Donación con usufructo y pacto de reversión}

\section{El problema: ¿rige el pacto de reversión después de 10 años?}

En la práctica es frecuente que quien dona un inmueble se reserve el usufructo y coloque un pacto de reversión para quedarse tranquilo. La discusión surge a partir de la postura según la cual el pacto de reversión no sería válido más allá de los 10 años. Desde el punto de vista registral, si el pacto de reversión no se levanta, sigue existiendo en el Registro de la Propiedad Inmueble.

\begin{note}
Las normas siguientes se transcriben en el material para facilitar el estudio; la docente se refiere al instituto sin indicar el número de artículo.

\textbf{Art. 1566 del Código Civil y Comercial de la Nación (pacto de reversión):} «En la donación se puede convenir la reversión de las cosas donadas, sujetando el contrato a la condición resolutoria de que el donatario, o el donatario, su cónyuge y sus descendientes, o el donatario sin hijos, fallezcan antes que el donante. Esta cláusula debe ser expresa y sólo puede estipularse en favor del donante. Si se la incluye en favor de él y de sus herederos o de terceros, sólo vale respecto de aquél. Si la reversión se ha pactado para el caso de muerte del donatario sin hijos, la existencia de éstos en el momento del deceso de su padre extingue el derecho del donante, que no renace aunque éste les sobreviva.»
% TODO: contenido incompleto o ambiguo en las notas originales (la redacción citada del art. 1566 se transcribe tal como figura en el apunte)

\textbf{Art. 1568 del Código Civil y Comercial de la Nación:} «La conformidad del donante para la enajenación de las cosas donadas importa la renuncia del derecho de reversión.» Esta norma guarda relación con la «renuncia a la reversión» que se menciona como requisito registral.
\end{note}

\subsection{Fundamento de la postura que limita el pacto a 10 años}

La postura que limita el pacto de reversión a 10 años lo asimila con la \textbf{indivisión forzosa} del Código Civil y Comercial. Los argumentos son los siguientes:

\begin{itemize}
    \item La indivisión forzosa también se establece por un plazo de 10 años, y la prescripción corta es igualmente de 10 años. Hay una idea de plazo: los derechos no pueden ser \textit{in aeternum}.
    \item Un pacto de reversión sin límite impediría un derecho constitucional como el derecho de propiedad. Si se dona y se transmite, la cláusula de reversión en algún momento debe tener un fin, pues de lo contrario el donatario nunca podría consolidar su titularidad sobre el inmueble. Se desdibuja el derecho de propiedad: sería «te lo doy, pero no te lo doy».
\end{itemize}

\subsection{Posición práctica}

La posición práctica sostenida en la clase es que, para afirmar que el pacto no subsiste más allá de los 10 años, es necesario \textbf{plantearlo judicialmente}, porque el Registro no levanta la inscripción si no hay renuncia a la reversión. Se requiere, entonces, una sentencia judicial de baja del pacto.

\begin{important}
Aunque pasen 10 años, el pacto de reversión sigue inscripto en el Registro de la Propiedad Inmueble. Para darlo de baja hace falta una \textbf{sentencia judicial} o una \textbf{renuncia voluntaria} de quien se reservó el derecho (por ejemplo, surgida de un acuerdo de mediación). Quien sostenga que caduca a los 10 años debe fundamentarlo en derecho y plantearlo judicialmente.
\end{important}

\section{Cómo se resuelve en sede registral}

Si transcurridos los 10 años el donatario quiere vender (por ejemplo, en el año 11), el Registro exigirá una de dos soluciones. Son las únicas dos formas de escriturar:

\begin{enumerate}
    \item la \textbf{renuncia previa} a la cláusula de reversión; o
    \item que el \textbf{adquirente declare que conoce y acepta} la cláusula.
\end{enumerate}

Quien quiera realizar este tipo de donación corre el riesgo de que el pacto de reversión genere un problema pasados los 10 años, salvo que un juez le dé la razón en el ámbito judicial o que surja de un acuerdo de mediación la renuncia voluntaria de quien se reservó el derecho. La existencia del problema se considera previsible.

\section{Derecho real: título, modo y registración}

\begin{definition}{Sistema argentino de adquisición de derechos reales}
El derecho real requiere \textbf{título} (causa fuente), \textbf{modo} (tradición posesoria) e \textbf{inscripción registral}, esta última con efectos \textbf{publicitarios}.
\end{definition}

En otros sistemas, como el sistema Torrens constitutivo, la inscripción perfecciona el derecho real. En el sistema argentino, en cambio, el derecho real se perfecciona en sede extrarregistral; la sede registral sirve para la publicidad \textit{erga omnes}, es decir, para quienes están fuera de la contratación. El contrato vale entre las partes aunque no esté inscripto, pero no será oponible a terceros si falta o adolece de registración.

\section{Caso práctico: planificación sucesoria y reversión}

Se relata un caso de planificación sucesoria de personas que no tienen una empresa familiar (o no les interesa como tal), son administradores de consorcio y poseen un patrimonio en inmuebles. Apareció una eventual nuera, que todavía no era nuera, y la titular del patrimonio manifestó no tener afinidad con ella; la planificación se organizó en consecuencia.

El problema de la reversión se presenta cuando se dona a alguien que está casado y que fallece antes que el donante. La cláusula opera de modo automático cuando el donatario muere antes que el donante. El fundamento que invoca el donante es que el cónyuge del donatario no participó en la generación del patrimonio, mientras que el donante, que sí lo generó, vería que su patrimonio pasa a quien no formó parte de su creación: el derecho sucesorio prevé que el cónyuge se quede con la mitad (la otra mitad iría a los padres si el donatario no tiene descendientes). Para evitarlo se pacta la reversión.

Cuando hay nietos, la cláusula de reversión suele perder sentido, porque existe continuidad del hijo o de la hija. Es una constante en las situaciones de planificación patrimonial sucesoria, y en ocasiones, ante la existencia de nietos, el donante renuncia a la reversión y se inscribe únicamente la renuncia, porque la familia ya se consolidó.

\begin{important}
La cláusula de reversión suele pactarse para evitar que, si el donatario fallece antes que el donante, el patrimonio llegue al cónyuge del donatario (que no participó en generarlo). Cuando nacen los nietos, la cláusula pierde sentido para el donante y es frecuente que renuncie a ella: se inscribe sólo la renuncia a la reversión.
\end{important}

\section{El factor humano: familias nuevas y sistemas «clánicos»}

El material señala que las nuevas configuraciones familiares se reflejan en la práctica de la planificación patrimonial. Se sostiene que no es posible querer a alguien que se conoce desde hace pocos años como si fuera un familiar, en el sentido de «darlo todo»; hacerlo exigiría una generosidad muy grande.

Se describe además el \textbf{sistema familiar «clánico»} (de \textit{clan}): es aquel que no se abre, y toda persona externa al clan es una extraña. Quien viene de afuera nunca llega a pertenecer, aunque aporte, mantenga o cuide a los miembros. Se afirma que esto no es ni bueno ni malo: es el parámetro del sistema familiar.

Esto se refleja en la práctica. En un matrimonio con tres hijos varones, dos nueras pueden estar fuera del clan, mientras que una puede estar en sintonía por haber sido criada con los mismos mandatos y principios. En ese caso hay afinidad y confianza, y a ese hijo se le proveerá de más. Incluso si la pareja se rompe, esa nuera puede seguir perteneciendo al sistema.

\section{Cuadro Cornell del capítulo}

\begin{longtable}{
    >{\raggedright\arraybackslash}p{0.27\textwidth}
    >{\raggedright\arraybackslash}p{0.65\textwidth}
}
\toprule
\textbf{Preguntas / palabras clave} & \textbf{Notas / respuestas} \\
\midrule
\endfirsthead
\toprule
\textbf{Preguntas / palabras clave} & \textbf{Notas / respuestas} \\
\midrule
\endhead
\bottomrule
\endfoot

\textbf{¿Caduca el pacto de reversión a los 10 años?} &
Es una discusión. La postura que lo limita lo asimila a la indivisión forzosa y a la prescripción corta (10 años) y se funda en que no se puede vaciar el derecho de propiedad (art.~17 CN). En la práctica, el pacto sigue inscripto en el Registro: se necesita sentencia judicial o renuncia voluntaria del donante (por ejemplo, vía mediación). \\[0.5em]

\textbf{¿Cómo se vende un inmueble con pacto de reversión inscripto después de 10 años?} &
Sólo hay dos formas de escriturar: previa renuncia a la cláusula de reversión, o que el adquirente declare que la conoce y la acepta. \\[0.5em]

\textbf{¿Qué son título, modo y registración?} &
Título es la causa fuente; modo es la tradición posesoria; la inscripción registral tiene efectos publicitarios (oponibilidad a terceros). En el sistema argentino el derecho real se perfecciona fuera del registro, a diferencia de sistemas constitutivos como el Torrens. \\[0.5em]

\textbf{¿Para qué se pacta la reversión y por qué pierde sentido con nietos?} &
Para evitar que el patrimonio donado pase al cónyuge del donatario si éste muere antes que el donante. Con la llegada de los nietos se consolida la familia y es frecuente que el donante renuncie a la reversión (se inscribe sólo la renuncia). \\[0.5em]

\textbf{¿Qué es un sistema familiar «clánico»?} &
Un sistema donde quien viene de afuera es siempre un extraño, aunque aporte y cuide. El profesional debe conocer los vínculos humanos porque se reflejan en la planificación patrimonial y sucesoria. \\

\midrule
\multicolumn{2}{p{0.92\textwidth}}{
\textbf{Resumen:} El pacto de reversión protege al donante frente al orden de los fallecimientos, pero su vigencia más allá de los 10 años es discutible. En la práctica, el Registro mantiene la inscripción hasta que haya renuncia o sentencia. El sistema argentino exige título, modo e inscripción registral con efectos publicitarios, y la planificación sucesoria debe considerar los vínculos familiares.
} \\
\end{longtable}

%==================================================================
% CAPÍTULO 3
%==================================================================
\chapter[Observabilidad de las donaciones de inmuebles]{Donaciones de inmuebles: antes de 2015, entre 2015 y 2020, y después de 2020}

\section{Planteo: el «limbo» legislativo}

El régimen de las donaciones de inmuebles presenta una situación de incertidumbre vinculada con los cambios normativos. Desde el punto de vista práctico, lo que importa es qué donaciones son observables y desde cuándo. En todos los casos hay que mirar \textbf{título por título} y calcular cuánto tiempo transcurrió.

\section{Donaciones anteriores a 2015}

Las donaciones anteriores a 2015 de padres a hijos no presentan inconvenientes. Respecto de las donaciones a terceros, la acción es de reducción de la herencia y, por el tiempo transcurrido, en muchos casos ya está prescripta.

\begin{formula}{Plazos de la acción en donaciones anteriores a 2015 (a terceros)}
La acción prescribe a los \textbf{10 años de la muerte del donante}, o a los \textbf{20 años desde el contrato de donación}.
\end{formula}

Lo relevante es la fecha de muerte del donante y no la del instrumento, que es la que normalmente circula. Por ello había que buscar la partida de defunción del donante. En ocasiones se abrían juicios sucesorios para demostrar que no había sucesores, pero eso no bastaba; lo único que subsanaba era el plazo.

El plazo de 20 años se vincula con un trabajo de ingreso a la Academia del Notariado del escribano Francisco Cerávolo, quien propuso la \textbf{prescripción adquisitiva por donación}. El razonamiento es que, si el Código Civil y Comercial ofrece la prescripción adquisitiva a quien no tenía título (causa fuente), con más razón corresponde permitirla cuando el donante, con voluntad propia, quiso dejar bien las cosas, sin que se forzara un contrato ni se mintiera diciendo que era una venta. La clase comparte este criterio.

\section{El Código de 2015: la parte luminosa y la negativa}

Con el Código de 2015 se buscó terminar con la incertidumbre sobre la fecha de muerte del donante. El Código computa los 10 años desde el acto de la \textbf{escritura de donación} o desde la \textbf{toma de posesión del donatario}. Puede haberse realizado un acta de entrega de posesión del inmueble (el modo, la \textit{traditio}) sin haberse celebrado el contrato de donación; si se logra demostrar la posesión como donatario, el plazo de 10 años comienza a correr desde ese momento. Esa es la parte luminosa de la norma.

La parte negativa es que extendió el régimen a todas las donaciones, incluso las de padres a hijos. Esto se considera violatorio del artículo 17 de la Constitución Nacional, porque el padre no podía planificar ni donar a sus propios hijos, restringiéndose el derecho constitucional de propiedad.

\begin{note}
\textbf{Constitución Nacional, art. 17 (mencionado en la clase):} «La propiedad es inviolable, y ningún habitante de la Nación puede ser privado de ella, sino en virtud de sentencia fundada en ley. La expropiación por causa de utilidad pública, debe ser calificada por ley y previamente indemnizada.» (...) Se reproduce el tramo que hace al derecho de propiedad, invocado como fundamento constitucional.
\end{note}

\section{La reforma de 2020 y la discusión actual}

En el año 2020, mediante el trabajo de las comisiones del Colegio de Escribanos, se revirtió legislativamente esta situación y se modificó el Código Civil: las donaciones de padres a hijos no tienen ningún tipo de observabilidad.

La norma es amplia: establece que se puede donar a cualquier persona sin observabilidad, con foco en que el instrumento pueda circular. Por ejemplo, si se dona a una institución como la Cruz Roja o la Fundación Favaloro, el título podría venderse y seguir circulando sin ser observable. Sin embargo, el notariado siguió discutiendo el punto y entiende que las únicas donaciones no observables son las de padres a hijos; las demás hay que analizarlas. Existen conclusiones de la Convención Notarial de 2024 en el Colegio de Escribanos donde se discutió nuevamente el tema. Hay quienes no están de acuerdo con la norma del Código Civil y Comercial que permite donar a cualquier persona.

Cuando se presenta una donación con fecha de instrumento entre 2015 y 2020, todas deben analizarse una por una: no hay un principio general. Lo que vale es la fecha de la escritura de donación o la toma de posesión, si se la puede acreditar. Todas esas donaciones serían observables.

En las restantes donaciones no corresponde mirar todo, sino \textbf{cuánto influye la donación en el patrimonio del donante}. Si entra en la porción disponible y no agrede la legítima de los herederos, el donatario no tendrá ninguna acción de reducción. La donación se computa en la cuenta de los bienes relictos dentro del sucesorio.

\subsection{Cuadro comparativo por época}

\begin{center}
\small
\begin{tabularx}{\textwidth}{
    >{\raggedright\arraybackslash}p{0.14\textwidth}
    >{\raggedright\arraybackslash}X
    >{\raggedright\arraybackslash}X
}
\toprule
\textbf{Período} & \textbf{Criterio / plazos} & \textbf{Observaciones} \\
\midrule
Antes de 2015 &
Donaciones de padres a hijos: sin inconvenientes. Con terceros: la acción prescribe a los 10 años desde la muerte del donante, o 20 años desde el contrato de donación. &
Se necesita la fecha de fallecimiento del donante, no la del instrumento. Cerávolo propuso la prescripción adquisitiva por donación. \\
\addlinespace
2015 a 2020 &
Los 10 años se computan desde la escritura de donación o desde la toma de posesión del donatario (si se la puede acreditar). La regla alcanzó a todas las donaciones, incluso las de padres a hijos. &
Se miran una por una, no hay un principio general. Se la considera violatoria del art.~17 de la Constitución Nacional. \\
\addlinespace
Desde 2020 &
Reforma legislativa lograda por las comisiones del Colegio de Escribanos: las donaciones de padres a hijos no tienen observabilidad. &
La norma es amplia (se puede donar a cualquiera), pero el notariado entiende que sólo las de padres a hijos no son observables. Ver conclusiones de la Convención Notarial 2024. \\
\bottomrule
\end{tabularx}
\end{center}

\begin{important}
\textbf{Regla práctica:} ante una donación de inmueble, mirar siempre título por título: fecha de la escritura, fecha de toma de posesión (si se puede acreditar), fecha de fallecimiento del donante (en las anteriores a 2015) y si la donación entra o no en la porción disponible sin afectar la legítima de los herederos.
\end{important}

\section{Defender una postura con fundamento}

Respecto de la discusión sobre el plazo del pacto de reversión, se enfatiza que el estudiante no debe quedar bien con nadie: debe saber \textbf{defender una postura} y justificarla, porque eso es lo que hará en la vida profesional. La postura que limita el pacto a 10 años puede ser comprendida y también refutada: el donatario tiene la expectativa de ejercer su derecho de propiedad libremente y sin restricciones (art.~17 de la Constitución Nacional, bases constitucionales del derecho privado). Desde ese argumento puede justificarse que el donatario no tiene expectativa de consolidar el bien si no se pone un límite a la cláusula de reversión: «te lo regalo, pero no te lo regalo».

Lo determinante es desde dónde se justifica: si se justifica en derecho y se justifica bien, no hay inconveniente.

\section{Cuadro Cornell del capítulo}

\begin{longtable}{
    >{\raggedright\arraybackslash}p{0.27\textwidth}
    >{\raggedright\arraybackslash}p{0.65\textwidth}
}
\toprule
\textbf{Preguntas / palabras clave} & \textbf{Notas / respuestas} \\
\midrule
\endfirsthead
\toprule
\textbf{Preguntas / palabras clave} & \textbf{Notas / respuestas} \\
\midrule
\endhead
\bottomrule
\endfoot

\textbf{¿Qué donaciones son observables y desde cuándo?} &
Antes de 2015: 10 años desde la muerte del donante o 20 desde el contrato (terceros). 2015--2020: 10 años desde la escritura o la toma de posesión, y alcanzaba incluso a padres a hijos (inconstitucional según la docente). Desde 2020: las de padres a hijos no son observables; las demás se miran una por una. \\[0.5em]

\textbf{¿Qué fecha importa en las donaciones anteriores a 2015?} &
La fecha de muerte del donante (partida de defunción) y no la del instrumento. Sólo el plazo subsana el problema. \\[0.5em]

\textbf{¿Qué propuso Cerávolo?} &
La prescripción adquisitiva por donación: si el Código ofrece prescripción adquisitiva a quien no tiene título, con más razón corresponde permitirla cuando hay donación voluntaria sin simulación de venta. \\[0.5em]

\textbf{¿Cuándo no hay acción de reducción contra el donatario?} &
Cuando la donación entra en la porción disponible y no agrede la legítima de los herederos; se computa en los bienes relictos dentro del sucesorio. \\[0.5em]

\textbf{¿Qué se espera del estudiante frente a posturas opuestas?} &
Que sostenga una postura y la justifique en derecho (Código y Constitución, art.~17 CN), sin buscar quedar bien con nadie. \\

\midrule
\multicolumn{2}{p{0.92\textwidth}}{
\textbf{Resumen:} La observabilidad de las donaciones de inmuebles depende de la época: antes de 2015, entre 2015 y 2020 y después de 2020. Siempre se debe analizar título por título (fecha de escritura, toma de posesión, fallecimiento del donante) y verificar si la donación entra en la porción disponible sin afectar la legítima. El criterio general es fundamentar jurídicamente la postura que se adopta.
} \\
\end{longtable}

%==================================================================
% CAPÍTULO 4
%==================================================================
\chapter[Tokenización inmobiliaria: concepto y valor]{Tokenización inmobiliaria: concepto, valor y diferencias con el fideicomiso}

\section{Qué es la tokenización}

Dentro de una red de \textit{blockchain} se genera un \textbf{token digital} que representa un activo que existe en la vida real y al cual está siempre vinculado. El token (por ejemplo, un NFT) está siempre vinculado con la propiedad o el activo representado. El estudio se centra en inmuebles, pero la tokenización también puede vincularse con bienes muebles.

La tokenización permite realizar \textbf{microtransacciones}: se pueden comprar tokens que representan fracciones mínimas del inmueble.

\begin{important}
Comprar un token \textbf{no} es adquirir el derecho real de dominio sobre el inmueble. Se adquiere un \textbf{derecho personal} y un derecho de propiedad sobre el token, un bien intangible que integra el patrimonio digital y está inescindiblemente vinculado a un activo del mundo real.
\end{important}

Lo esencial es escindir el derecho de propiedad del inmueble, que no cambia, porque seguirá a nombre de quien esté tokenizando el inmueble, de la titularidad del token. El token es un bien intangible que forma parte de lo que se denomina \textbf{patrimonio digital} de las personas.

\section{Qué se compra jurídicamente}

El token representa digitalmente la parte del inmueble, pero a los efectos jurídicos no se compra el inmueble. Se adquiere un \textbf{derecho personal de propiedad} sobre el token. Se advierte que debe hablarse de «propiedad» y no de «dominio»: el dominio, derecho real, no es lo mismo que la propiedad, que es un género abarcativo. Puede haber propiedad sobre un derecho personal, en este caso sobre el token, que es un intangible; son ceros y unos, un \textit{bit}. Representativamente, el token está vinculado con un bien del mundo real, por lo que se lo denomina \textit{real world asset} (RWA, activo del mundo real). Ese vínculo busca que la gente confíe en dónde invierte y que no se trate de un fideicomiso financiero y nada más.

Una asimilación útil, planteada en clase, es compararlo con comprar el derecho sobre una patente: se trata de un intangible.

\begin{note}
Otras jurisdicciones con un sistema de derechos reales distinto del argentino, como Dubái o Arabia Saudita, consideran que al comprar el token se compra la parte del inmueble, y eso es susceptible de registración. En el sistema argentino de título, modo y registración se necesita \textbf{escritura pública} para transmitir el dominio.
\end{note}

\section{No es un esquema Ponzi: el valor atado al inmueble}

Se sostiene que la tokenización no es un esquema Ponzi, porque está atada al valor que le da la propiedad. Si el valor de la propiedad sube, el valor del token acompaña. Ejemplo planteado: una microinversión de 100 pasa a valer 120 y, si la propiedad sigue subiendo, 150, lo que implica un 50\,\% ganado sobre la inversión originaria.

\begin{example}{Evolución del valor del token}
Se invierten 100. Si la propiedad sube, el token pasa a valer 120. Si la propiedad sigue subiendo, el token pasa a valer 150 al momento de venderlo: se obtuvo una ganancia del 50\,\% sobre la inversión original.
\end{example}

\subsection{Blockchain como «tecnología del valor»}

Por esto se conoce a \textit{blockchain} como la «tecnología del valor». Las figuras de la Web~3, como tecnologías del valor, permiten al sujeto mantener el valor económico de su inversión, o mejorarlo. Se vincula con un pensamiento filosófico inglés liberal extremo, según el cual «los Estados nos hacen perder dinero»: se pueden generar espacios colaborativos, compartidos y descentralizados en los que el patrimonio mantenga su valor sin verse depreciado por malos manejos económicos o por inflación.

La mayoría de las aplicaciones corren por \textbf{Ethereum}, porque es la única red de \textit{blockchain} que permite lo conocido como \textit{layers} (capas), en las cuales se permite la redacción informática de contratos inteligentes (\textit{smart contracts}).
% TODO: contenido incompleto o ambiguo en las notas originales (fragmento incompleto en la grabación)

\section{Smart contracts, white paper y reserva de valor}

Los \textit{smart contracts} permiten programar reglas. Por ejemplo, que el token sólo pueda venderse a partir de que el inmueble tenga construido su quinto piso, siendo un edificio de seis pisos: hasta entonces el inversor permanece en la inversión y no puede salir. Todas esas reglas están redactadas en el \textbf{white paper} del proyecto.
% TODO: contenido incompleto o ambiguo en las notas originales (la definición de white paper no figura en el apunte; fragmento incompleto en la grabación)

Una de las ventajas es que el valor pagado por la unidad resulta considerablemente menor que el desembolso de comprarla completa en el mercado en un único pago. El valor del token es distinto si se compra cuando el desarrollador empieza la demolición de la casona adquirida, que cuando el edificio ya está listo para entregar la posesión. Ese es el lugar de la \textbf{reserva de valor} del inversor.

\section{Qué cambia la tokenización respecto del fideicomiso}

En el contrato de fideicomiso, la transmisión (ya sea por adhesión al fideicomiso o porque el fiduciario firmó un boleto de compraventa) recae siempre sobre el \textbf{100\,\% de la futura unidad funcional}. La tokenización rompe este esquema, porque permite comprar una \textbf{fracción} de la unidad. Un inversor joven, sin ahorros suficientes para la unidad completa, puede comprar por ejemplo un 25\,\% en una primera inversión y, tras la salida de ese proyecto, acceder luego a una unidad mayor; en un segundo proyecto de tokenización podría adquirir la vivienda.

Uno de los objetivos es permitir un acceso a la vivienda de manera más amigable. En la economía argentina, la gente joven, si no recibe ayuda de su familia, no logra acceder a la vivienda, y los créditos no alcanzan para satisfacer la demanda habitacional. Se buscan, por ello, formas de financiación y de cultura financiera.
% TODO: contenido incompleto o ambiguo en las notas originales (fragmento incompleto en la grabación)

Otras características planteadas:

\begin{itemize}
    \item En la plataforma se colocan únicamente los tokens efectivamente adquiridos.
    \item Existe una diferencia entre el \textbf{criptoactivo como género} y los \textbf{tokens como especie}: los tokens participan de un proyecto de tokenización específico.
    \item El token puede usarse como \textbf{instrumento de pago}, cuando hay acuerdos con otros socios del mismo ecosistema que aceptan los tokens del proyecto.
    \item Funciona como \textbf{reserva de valor}, porque permite capitalizarse, y constituye una herramienta de \textbf{financiación privada}, que no depende de una cartera de préstamos hipotecarios ni de políticas públicas del Estado, sino de la iniciativa de los particulares.
\end{itemize}

\subsection{Cuadro comparativo: fideicomiso inmobiliario vs. tokenización}

\begin{center}
\small
\begin{tabularx}{\textwidth}{
    >{\raggedright\arraybackslash}p{0.17\textwidth}
    >{\raggedright\arraybackslash}X
    >{\raggedright\arraybackslash}X
}
\toprule
\textbf{Aspecto} & \textbf{Fideicomiso inmobiliario} & \textbf{Tokenización} \\
\midrule
Naturaleza &
Contrato con diferentes cláusulas; se asemeja a la compra de un inmueble en cuotas. &
Digitalización del fideicomiso; inversión en tokens vinculados a un activo real (RWA). \\
\addlinespace
Objeto de la transmisión &
Siempre sobre el 100\,\% de la futura unidad funcional. &
Fracciones mínimas: microtransacciones y microinversiones. \\
\addlinespace
Obligación de adquirir la unidad &
Sí: aportes durante la obra y eventuales diferencias por mayores costos. &
No: se puede capitalizar, vender en el mercado secundario o quemar tokens para adjudicarse otra unidad. \\
\addlinespace
Mirada del participante &
Sobre la unidad en sí (deseo de tener casa propia). &
Sobre la inversión (tecnología del valor). \\
\addlinespace
Acceso &
Adhesión al fideicomiso o boleto de compraventa con el fiduciario. &
Plataforma digital: \textit{onboarding}, controles de lavado, pago con tarjeta y contrato de mandato de token. \\
\addlinespace
Derecho que se adquiere &
Derecho personal hasta la transmisión del dominio. &
Derecho personal de propiedad sobre el token (bien intangible, patrimonio digital). \\
\addlinespace
Causa fuente del dominio &
Contrato de fideicomiso / compraventa. &
El proceso de tokenización documentado en el \textit{white paper}. \\
\addlinespace
Mercado secundario &
No previsto en el esquema. &
Sí, según lo que disponga el \textit{white paper}. \\
\bottomrule
\end{tabularx}
\end{center}

\section{Cuadro Cornell del capítulo}

\begin{longtable}{
    >{\raggedright\arraybackslash}p{0.27\textwidth}
    >{\raggedright\arraybackslash}p{0.65\textwidth}
}
\toprule
\textbf{Preguntas / palabras clave} & \textbf{Notas / respuestas} \\
\midrule
\endfirsthead
\toprule
\textbf{Preguntas / palabras clave} & \textbf{Notas / respuestas} \\
\midrule
\endhead
\bottomrule
\endfoot

\textbf{¿Qué es la tokenización?} &
La representación digital, en una red \textit{blockchain}, de un activo del mundo real (RWA) mediante tokens que permiten comprar fracciones mínimas (microtransacciones). \\[0.5em]

\textbf{¿Qué se adquiere al comprar un token?} &
Un derecho personal de propiedad sobre el token (bien intangible, parte del patrimonio digital), no el dominio del inmueble. La propiedad sigue a nombre de quien tokeniza. Propiedad (género) no es lo mismo que dominio (derecho real). En otras jurisdicciones (Dubái, Arabia Saudita) se considera que se compra la parte del inmueble. \\[0.5em]

\textbf{¿Cómo se genera valor con un token?} &
El valor del token está atado al valor del inmueble; si éste sube, el token también. \textit{Blockchain} se conoce como «tecnología del valor» porque permite mantener o mejorar el valor de la inversión. \\[0.5em]

\textbf{¿Qué es el white paper y qué son los smart contracts?} &
El \textit{white paper} contiene las reglas del proyecto (cantidad de tokens, fraccionamiento, condiciones de venta, obligaciones del desarrollista). Los \textit{smart contracts} son contratos programados, típicamente sobre Ethereum, que ejecutan esas reglas (por ejemplo, permitir la venta recién cuando se construye el quinto piso). \\[0.5em]

\textbf{¿Qué cambia la tokenización respecto del fideicomiso?} &
El fideicomiso transmite siempre el 100\,\% de la futura unidad; la tokenización permite comprar fracciones, no obliga a adquirir la unidad, admite mercado secundario y se mira como inversión. Facilita un acceso más amigable a la vivienda y constituye financiación privada. \\

\midrule
\multicolumn{2}{p{0.92\textwidth}}{
\textbf{Resumen:} Un token es la representación digital de un activo real (RWA) en una red \textit{blockchain}. Comprarlo no transmite el dominio sino un derecho personal de propiedad sobre un bien intangible. Su valor está atado al inmueble, sus reglas se fijan en el \textit{white paper} y se ejecutan mediante \textit{smart contracts}. Frente al fideicomiso, permite fracciones, no obliga a adquirir la unidad y habilita un mercado secundario.
} \\
\end{longtable}

%==================================================================
% CAPÍTULO 5
%==================================================================
\chapter[Regulación: sandbox normativo, CNV y PSAV]{Regulación: sandbox normativo, CNV y proveedores de servicios de activos virtuales}

\section{El sandbox normativo}

\begin{definition}{Sandbox normativo}
Es el conjunto de resoluciones que genera el Estado para hacer amigable el trabajo de los privados y permitir nuevos negocios, con una regulación menor pero contenida.
\end{definition}

En lo técnico, \textit{sandbox} es una forma de aislar de internet o de elementos exteriores para dar más seguridad al activo. Pero en el plano normativo, el \textit{sandbox} es el «arenero»: la metáfora es la del tobogán que desemboca en arena para no golpearse contra el cemento. Un \textit{sandbox} normativo es un conjunto de disposiciones legales que permite a las partes trabajar menos reguladas, pero contenidas: que no puedan hacer daño a terceros, que no puedan fugarse con los capitales recibidos y que no ingrese dinero del lavado. Se respetan los marcos generales de la ley, pero se permite trabajar de manera que el negocio sea viable.

\begin{note}
\textbf{Referencia normativa: Resoluciones Generales de la Comisión Nacional de Valores (CNV).} Se citan la \textbf{RG 1137}, que inicia el régimen de tokenización y crea el registro de los proveedores de servicios de activos virtuales, y la \textbf{RG 1150}, que amplía el régimen de tokenización de la RG 1137 y prorroga el \textit{sandbox} normativo hasta 2027 («un año más de sandbox normativo»). La CNV empieza a dictar estas resoluciones a partir de 2025. No se dicta el texto de las normas.
% TODO: contenido incompleto o ambiguo en las notas originales (se menciona «la última» como «1057» de manera dubitativa; se promete un cuadro con las resoluciones)
\end{note}

La idea del \textit{sandbox} es invitar a incubar y trabajar «en laboratorio» dentro de un contexto marco.

\section{Los PSAV y los riesgos que se quieren evitar}

El contexto marco nace con la creación de los \textbf{proveedores de servicios de activos virtuales (PSAV)}. Estos proveedores:

\begin{itemize}
    \item deben estar \textbf{registrados};
    \item deben dejar una \textbf{caución de fondos} para asegurar que no se fugarán con la inversión que ingresó a sus plataformas;
    \item deben cumplir la normativa de prevención de lavado de activos, financiamiento del terrorismo y proliferación de armas de destrucción masiva; y
    \item deben cumplir con ARCA.
\end{itemize}
% TODO: contenido incompleto o ambiguo en las notas originales (en la transcripción original la sigla figura como «PCAB»; se unifica como PSAV por el contexto)

La tokenización se desarrolla en un entorno muy regulado porque existe el riesgo de ingresar a un \textbf{esquema Ponzi} sin saberlo, o a una plataforma que muestra un \textit{dashboard} (pantalla operativa) pero detrás de la cual hay un atacante que, al transferirse el dinero, se lo lleva. Sin el \textit{sandbox} regulatorio habría más expedientes judiciales por casos Ponzi o por hackeo.

\section{El white paper y el rol del profesional}

Cada proyecto debe tener su \textit{white paper}, y el profesional debe aprender a leerlos, porque es lo que llevará el cliente cuando diga que quiere invertir y solicite asesoramiento remunerado.

Como analogía, si un cliente quiere comprar una unidad «en pozo» regida por un fideicomiso, lo primero que se revisa es el contrato (que es abierto, por lo que las cláusulas pueden ser infinitas) y si el fideicomiso está inscripto.

\begin{important}
\textbf{Primer requisito ante un cliente que quiere comprar tokens:} verificar que el proveedor de servicios de activos virtuales esté inscripto. Si no lo está, es una estafa o un esquema Ponzi.
\end{important}

\section{Organismos intervinientes}

El fideicomiso, al menos en la Ciudad de Buenos Aires, se inscribe primero en la \textbf{Inspección General de Justicia} y luego, por la parte administrativa, en el \textbf{Gobierno de la Ciudad de Buenos Aires}; además, el contrato de fideicomiso se inscribe en el \textbf{Registro de la Propiedad Inmueble}. La tokenización, por ahora, no tiene inscripción del proyecto: tiene un \textit{white paper}. Su regulación consiste en que el PSAV que interviene en la plataforma debe estar inscripto en el régimen que crea la RG 1137. Cuando se habla de tokenización se habla únicamente de la \textbf{Comisión Nacional de Valores}, por ahora, hasta que el proyecto se termine.

\begin{center}
\small
\begin{tabularx}{\textwidth}{
    >{\raggedright\arraybackslash}p{0.22\textwidth}
    >{\raggedright\arraybackslash}X
    >{\raggedright\arraybackslash}X
}
\toprule
\textbf{Figura} & \textbf{Organismos / registros} & \textbf{Cuestión clave} \\
\midrule
Fideicomiso inmobiliario (CABA) &
Inspección General de Justicia, luego Gobierno de la Ciudad (cuestiones administrativas) y Registro de la Propiedad Inmueble (inscripción del contrato). &
Inscripción múltiple del fideicomiso. \\
\addlinespace
Tokenización &
Comisión Nacional de Valores (RG 1137, ampliada por la RG 1150). &
No hay inscripción del proyecto, sino un \textit{white paper}; el PSAV debe estar inscripto. \\
\bottomrule
\end{tabularx}
\end{center}

\section{Cuadro Cornell del capítulo}

\begin{longtable}{
    >{\raggedright\arraybackslash}p{0.27\textwidth}
    >{\raggedright\arraybackslash}p{0.65\textwidth}
}
\toprule
\textbf{Preguntas / palabras clave} & \textbf{Notas / respuestas} \\
\midrule
\endfirsthead
\toprule
\textbf{Preguntas / palabras clave} & \textbf{Notas / respuestas} \\
\midrule
\endhead
\bottomrule
\endfoot

\textbf{¿Qué es el sandbox normativo?} &
Un paquete de resoluciones del Estado (CNV) que permite trabajar con menor regulación pero contenida: sin dañar a terceros, sin fuga de capitales y sin lavado. Es la metáfora del «arenero». Se prorrogó hasta 2027 (RG 1150, que amplía la RG 1137). \\[0.5em]

\textbf{¿Qué verificar primero si un cliente quiere comprar tokens?} &
Que el proveedor de servicios de activos virtuales (PSAV) esté inscripto en la CNV. Si no lo está, es una estafa o un esquema Ponzi. Los PSAV dejan una caución, previenen el lavado y cumplen con ARCA. \\[0.5em]

\textbf{¿Qué organismos intervienen en el fideicomiso y en la tokenización?} &
Fideicomiso inmobiliario: IGJ, Gobierno de la Ciudad y Registro de la Propiedad Inmueble. Tokenización: sólo CNV (por ahora); no hay inscripción del proyecto sino un \textit{white paper}. \\[0.5em]

\textbf{¿Por qué se regula con un sandbox?} &
Para evitar esquemas Ponzi y ataques en los que, detrás de una plataforma aparente, un atacante se lleva el dinero transferido, y para no inundar de expedientes judiciales. \\

\midrule
\multicolumn{2}{p{0.92\textwidth}}{
\textbf{Resumen:} El Estado regula la tokenización mediante un \textit{sandbox} normativo de la CNV (RG 1137 y RG 1150, prorrogado hasta 2027), que permite trabajar con una regulación menor pero contenida. Los PSAV deben estar inscriptos, dejar una caución y cumplir con la prevención de lavado y con ARCA. Mientras el fideicomiso se inscribe en IGJ, Gobierno de la Ciudad y Registro, la tokenización depende por ahora de la CNV y del \textit{white paper}.
} \\
\end{longtable}

%==================================================================
% CAPÍTULO 6
%==================================================================
\chapter[Ciclo de vida del token y causa fuente]{Ciclo de vida del token, mercado secundario y causa fuente del dominio}

\section{Minteo, transferencia y quema}

El token tiene tres momentos en su ciclo de vida.

\begin{enumerate}
    \item \textbf{Emisión o minteo (\textit{minting}):} es la generación de los tokens. En el sistema teórico se emiten tantos tokens como partes del inmueble se van a vender. En los entornos controlados que se aplican actualmente, se emiten en la medida en que hay adhesiones al fideicomiso.
    \item \textbf{Transferencia:} se valida que el titular efectivamente tenía la cantidad de tokens y que puede transmitirlos; aquí se encuentra el mercado secundario.
    \item \textbf{Rescate o quema (\textit{burning}):} es como el rescate de las acciones. Se queman tokens y, en función de la cantidad de tokens que se tiene, se accede al 100\,\% de una unidad (por ejemplo, el primero A del inmueble de determinada calle). Llega el momento en que el titular se va del proyecto con las ganancias y el capital aportado, o se adjudica los derechos a determinada unidad, o los vende a otro en el mercado secundario.
\end{enumerate}

\begin{center}
\small
\begin{tabularx}{\textwidth}{
    >{\raggedright\arraybackslash}p{0.30\textwidth}
    >{\raggedright\arraybackslash}X
}
\toprule
\textbf{Etapa} & \textbf{Qué ocurre} \\
\midrule
1. Emisión o minteo (\textit{minting}) &
Generación de los tokens. En el sistema teórico se emiten tantos tokens como partes se van a vender; en Argentina se emiten a medida que hay adhesiones. \\
\addlinespace
2. Transferencia &
Se valida que el titular tiene la cantidad de tokens y puede transmitirlos; mercado secundario. \\
\addlinespace
3. Rescate o quema (\textit{burning}) &
Se rescatan los tokens: el titular se va con ganancias y capital, o se adjudica una unidad, o los vende a otro. \\
\bottomrule
\end{tabularx}
\end{center}

\section{La causa fuente del derecho real cambia}

La normativa de derechos reales sigue siendo la del Código Civil y Comercial. Lo que cambia es la \textbf{causa fuente} del derecho real de dominio: no es la compraventa ni el contrato de fideicomiso; la causa fuente es el \textbf{proceso de tokenización}. Siguen rigiendo título, modo y registración.

\begin{important}
La causa fuente cambia, el régimen no: la transmisión del dominio sigue regida por el Código Civil y Comercial (título, modo, registración y escritura pública), pero el título ya no es una compraventa ni una adhesión al fideicomiso individual: es el proceso de tokenización documentado en el \textit{white paper}.
\end{important}

En la escritura no se dirá que una persona, en determinada fecha, suscribió una adhesión al contrato de fideicomiso o firmó un boleto de compraventa con el fiduciario. Se dirá que, en determinada fecha, suscribió cierta cantidad de tokens dentro del proyecto, en el marco del \textit{white paper}; que se asumieron determinadas obligaciones; que en otras fechas compró más tokens; y que todo ello generó el derecho, en función de una cláusula del \textit{white paper}, a adjudicarse la unidad funcional correspondiente, y que en cumplimiento de todo esto se transmite el dominio.

No hay una venta de un tercero ni del constructor: es necesario relatar todo el proceso, porque la causa surge de allí. No se compra poniendo dinero sobre la mesa, sino que se realizó un proceso para gestionarse los tokens y luego adjudicarlos. La situación es equivalente a haber firmado pequeñas adhesiones al fideicomiso hasta reunir toda la unidad. Esto es propio de la materia de derechos reales, porque se necesita la causa jurídica de la adquisición.

Si alguien compra todos los tokens de una unidad funcional y quiere materializar el bien (por ejemplo, para ir a vivir), igualmente debe cumplirse todo el procedimiento ordinario de derechos reales.

\section{Comparación práctica: aporte en el fideicomiso y compra de tokens}

En el \textbf{fideicomiso} se aporta al pozo a medida que avanza la construcción; puede haber imprevistos y exigirse aportes adicionales en determinada etapa, o diferencias por mayores costos. El adherente (o el comprador, si fue por boleto) asume obligaciones. Por eso es fundamental saber leer las cláusulas, lo mismo que se hará con el \textit{white paper}, porque puede ocurrir que en el \textit{white paper} se vuelque un contrato de fideicomiso bajo la modalidad de \textit{white paper} del token, y que se lo acepte en los términos y condiciones, constituyendo ésa la firma.

En la \textbf{tokenización}, la persona entra en una página, paga con tarjeta de crédito lo que quiere invertir y puede comprar tokens todos los meses, según sus posibilidades. Puede tener una idea estimada de cuánto representará lo que compra al final, pero no lo sabe exactamente; conoce cuánto tiene capitalizado en ese momento, cuánto representa y cuánto representa el valor de la unidad.

Hoy pueden existir inmuebles bajo ambos regímenes: una parte bajo el régimen del fideicomiso y otra parte del edificio vendida tokenizada. Hay constructores que dejan dos unidades para la tokenización y venden el resto por el sistema tradicional.

\section{Mercado secundario y tecnología del valor: el ejemplo de las fracciones}

\begin{example}{Compra de tokens en el mercado secundario}
A Carolina le faltan 30 tokens para quedarse con el primero A. Diez los tiene Matías, diez Tamara y diez Santi. Carolina deberá pagar un valor distinto por esos tokens. Para Matías, Tamara y Santi fue una muy buena inversión: pensaban ir de vacaciones a Brasil, pero como lograron vender tan bien los tokens, viajarán a Europa. Carolina se quedó con el primero A.
\end{example}

La diferencia es que \textbf{no hay obligación de comprar la unidad}. Comprarla resulta más inestable, porque depende de que otros quieran vender. Quien tiene tokens suficientes para el primero A, pero encuentra libre el segundo B, puede aplicarlos, quemarlos y adjudicarse el segundo B. Hay mucha libertad de maniobra: puede no querer un inmueble en ese lugar, haberse capitalizado en el proyecto, quemar los tokens, hacerse de ese valor y comprar de manera tradicional en otro lugar.

Por eso se habla de tecnología del valor. En el fideicomiso está el deseo de adquirir un inmueble o una casa propia y se mira siempre la unidad en sí; en la tokenización se mira la \textbf{inversión}, vinculada a un \textit{real world asset} (RWA).

\subsection{La analogía del álbum de figuritas}

El álbum de figuritas ayuda a visualizar el mecanismo: se necesita una figurita determinada para completar el álbum, así como se necesita el token determinado para acceder a la unidad.

\section{Real World Assets y oráculos}

Un \textit{Real World Asset} es, en el mundo de la tokenización, la vinculación del token con un bien que está \textit{off chain}, en el mundo real, y al cual el token queda necesariamente vinculado; de allí la inescindibilidad entre inmueble y token.

El mercado inmobiliario seguirá funcionando, según las conclusiones del congreso internacional de derecho inmobiliario y tecnología, porque alguien debe verificar que quien incorpora el inmueble a la plataforma es dueño del inmueble, que el inmueble no tenga embargos y que el dueño, que es quien desarrolla, no tenga inhibiciones. Lo mismo que sucede en el fideicomiso tradicional se continuará haciendo en el mundo virtual bajo la denominación de \textbf{oráculos}.

\begin{definition}{Oráculo}
Es una persona designada en el \textit{white paper}, una organización designada u otro programa informático designado para buscar información en la vida real y traerla al contrato inteligente. Cuando el contrato inteligente debe operar, solicita información al oráculo y, según la información que éste trae, ejecuta una acción u otra.
\end{definition}

Se habla de convergencia porque intervienen inteligencia artificial, \textit{blockchain}, \textit{software}, plataformas, ciberseguridad y experiencia de usuario. Todo se basa en la seguridad que se brinda al inversor: que no está comprando humo y que lo invertido está vinculado a la historia del inmueble que se construye. Se deberá verificar que el constructor cumpla y que se cumplan los plazos, salvo cuestiones climáticas o imponderables, que estarán en el \textit{white paper}, al igual que están en las cláusulas del contrato de fideicomiso, y todo respaldado por una red de \textit{blockchain}, con nodos y criptografía.

\section{Cuadro Cornell del capítulo}

\begin{longtable}{
    >{\raggedright\arraybackslash}p{0.27\textwidth}
    >{\raggedright\arraybackslash}p{0.65\textwidth}
}
\toprule
\textbf{Preguntas / palabras clave} & \textbf{Notas / respuestas} \\
\midrule
\endfirsthead
\toprule
\textbf{Preguntas / palabras clave} & \textbf{Notas / respuestas} \\
\midrule
\endhead
\bottomrule
\endfoot

\textbf{¿Cuál es el ciclo de vida del token?} &
Minteo (emisión; en Argentina a medida que hay adhesiones), transferencia (mercado secundario) y quema o \textit{burning} (rescate: adjudicarse la unidad, vender o irse con las ganancias). \\[0.5em]

\textbf{¿Qué cambia en el derecho real con la tokenización?} &
Cambia la causa fuente: ya no es una compraventa ni una adhesión al fideicomiso, sino el proceso de tokenización en el marco del \textit{white paper}. Siguen rigiendo título, modo, registración y escritura pública (Código Civil y Comercial). \\[0.5em]

\textbf{¿Hay obligación de adquirir la unidad?} &
No. A diferencia del fideicomiso, el titular puede capitalizarse, vender en el mercado secundario o quemar tokens para adjudicarse otra unidad. \\[0.5em]

\textbf{¿Qué es un Real World Asset?} &
La vinculación del token con un bien \textit{off chain}, del mundo real, de modo inescindible. \\[0.5em]

\textbf{¿Qué son los oráculos?} &
Personas, organizaciones o programas designados en el \textit{white paper} que traen al contrato inteligente información del mundo real (dominio, embargos, inhibiciones) para que decida cómo operar. \\

\midrule
\multicolumn{2}{p{0.92\textwidth}}{
\textbf{Resumen:} El token atraviesa minteo, transferencia y quema. Cambia la causa fuente del dominio (el proceso de tokenización documentado en el \textit{white paper}), pero no el régimen del Código Civil y Comercial, que exige título, modo, registración y escritura. No hay obligación de adquirir la unidad, y los oráculos permiten verificar en el mundo real la información que el contrato inteligente necesita.
} \\
\end{longtable}

%==================================================================
% CAPÍTULO 7
%==================================================================
\chapter[Blockchain aplicada y token de calidad]{Blockchain aplicada: nodos, hash, onboarding y token de calidad}

\section{Nodos y consenso}

Los tokens están creados y asegurados criptográficamente. Intervienen la \textbf{clave asimétrica} y el \textbf{hash} que se produce cuando se cierra cada bloque y se incorpora el \textit{time stamp} o sello de tiempo, lo que permite acreditar que determinado token pertenece a una persona y no a otra.

Los \textbf{nodos} son las computadoras que conforman la red y funcionan como un libro contable en tiempo real. Todos tienen información y acceso a todo. Hay tantas copias del libro como nodos en la red; todos tienen el mismo contenido y generan el \textbf{consenso}, que es la pauta de verdad. Como todos tienen lo mismo y están de acuerdo, se pueden validar las transacciones.

\section{El hash en un caso concreto}

El hash es la incógnita que debe resolverse dentro de la cadena de \textit{blockchain} para guardar la transacción. En el ejemplo, la transacción en la que Carolina le compra tokens a Matías queda validada por el hash, que representa la transferencia realizada. Esa transacción queda resguardada en el bloque, y el hash del bloque la documenta; es equivalente a que hubiera quedado documentada la cesión de derechos del fideicomiso ante el escribano.

\begin{formula}{Hash}
El hash es un código alfanumérico de \textbf{64 caracteres}, que se trabaja con un cálculo computacional para que los mineros logren la encriptación del bloque, la generación del bloque y la creación de la marca temporal (\textit{timestamp}, sello de tiempo) que valida la transacción.
\end{formula}

Estos contenidos se relacionan con la clave asimétrica (claves públicas y privadas, identificación del usuario) y con la firma digital: se avanza paso a paso desde la firma digital, luego \textit{blockchain} y finalmente la tokenización.

Cuando se adquiere un token, esa adquisición queda registrada en la red y no puede revertirse ni borrarse. Allí radica la \textbf{inmutabilidad} de la red de \textit{blockchain} y el \textbf{sello de tiempo} de la generación del bloque, que impide que se pueda quitar al titular el derecho a ese token.

\section{Pérdida de usuario y clave: diferencia con las criptomonedas}

En las criptomonedas, si se pierden el usuario y la clave, se pierde todo. En la tokenización no ocurre lo mismo, porque existe un paso previo: el \textbf{proceso de autenticación del usuario}. Al ingresar a la plataforma se pasan controles de lavado y se solicita el ingreso con datos personales; es probable que intervengan aplicaciones \textit{fintech} que validan la identidad mediante datos biométricos (mirar a la derecha o a la izquierda, fotografía con el documento). Todo ello forma parte del \textit{onboarding}.

\section{El rol del escribano en el onboarding}

Algunos proyectos de tokenización piden que el \textit{onboarding} lo haga el escribano. Allí la intervención notarial recupera relevancia: es quien da el acceso mediante la gestión de identidades, a través de una adhesión al fideicomiso realizada completamente en línea y de la certificación de firmas del sistema de firma a distancia que desarrollan los colegios. En el caso de la Capital, el sistema es el SDF 2.0: en una reunión virtual la persona firma el contrato de token (llamado \textbf{mandato de token}) y realiza toda su adhesión al fideicomiso; con esa adhesión se genera el token.

\section{Proyecto con convenio notarial y «token de calidad»}

Se menciona el proyecto de Pala Blockchain, con convenio marco firmado con el Colegio de Escribanos de la Ciudad de Buenos Aires. Se diferencia de las tokenizaciones tradicionales porque desarrolla un \textbf{token de calidad}, denominado «casa token», que es un receptor de información: dentro del token están el contrato de fideicomiso, la adhesión de la persona y los controles de identidad realizados respecto de la prevención de lavado de activos y el «conozca a su cliente», es decir, los controles que deben tener los PSAV.
% TODO: contenido incompleto o ambiguo en las notas originales (los nombres «Pala Blockchain» y «token de calidad / casa token» figuran así en la transcripción original; conviene verificarlos con el material de la cátedra)

El token contiene entonces el contrato de mandato, datos personales, datos del proyecto, datos del inmueble y datos de lo adquirido. Por eso, si se pierden usuario y clave, pueden ser restablecidos: la transacción tiene un sello de tiempo en \textit{blockchain}, que corre por la red de \textbf{Polygon}, otra red de \textit{blockchain} además de Ethereum. La transacción queda registrada en el bloque de Polygon, por lo que el desarrollador no podría desconocer posteriormente la compra del token. Esto también favorece la ciberseguridad, porque el nombre ya está registrado.

\section{Tokenización clásica y tokenización en Argentina}

La \textbf{tokenización clásica} consiste en emitir todos los tokens representados para que luego las personas los compren. Metro Futuro lo hace así, pero en España, y lo que tokeniza no es el inmueble sino las acciones de la sociedad propietaria del inmueble (la SL): se tokeniza el 100\,\% de las acciones, que representan lo único que tiene la sociedad, el inmueble.

En el caso de Pala Blockchain se admite que no se realiza una tokenización pura, es decir, que no se emiten todos los tokens para que sean suscriptos, sino que se los \textbf{va subiendo a medida que la persona se adhiere}, lo que genera menos riesgo. El token de calidad, mientras no esté lanzado como token, sigue siendo propiedad del desarrollista, como si éste los hubiera emitido y los tuviera. Si el desarrollista no emitió todos los tokens y el proyecto fracasa, seguirá siendo el propietario, del mismo modo que ocurre en el fideicomiso.

\section{Dilución del valor y fraccionamiento}

Se plantea la inquietud de si el desarrollista que conserva el token de calidad podría poner a la venta cantidades parciales y luego liberar más para obtener mayor inversión, como ocurre en el fideicomiso analógico cuando el desarrollador retiene unidades para venderlas a mayor valor.

La respuesta de la clase es que el desarrollista \textbf{no puede licuar el valor del token}: no le conviene, porque nadie invertiría. Tampoco puede dividir el token: la división ya se produjo, el token es uno solo. El token adquirido, registrado en \textit{blockchain}, no puede ser tocado ni subdividido por nadie. Luego el inversor sigue comprando la cantidad que tiene y decide en qué unidad aplicarla.

El nivel de fraccionamiento se fija en el \textit{white paper}: allí debe constar en cuántos tokens se dividirá el proyecto. En la plataforma, en cambio, se ven únicamente los tokens que efectivamente tienen dueño, los que se pusieron a la venta. No existen tokens vacíos disponibles como en otras plataformas: la persona debe pasar por el \textit{onboarding}, los controles de lavado y la suscripción del contrato de mandato, y compra su primer token. Si dos meses después vuelve a invertir, compra otro mandato de token y obtiene otro token, sin repetir el \textit{onboarding}.

Este entorno más controlado evita que se venda dos veces lo mismo. Según el modo en que lo explica Pala Blockchain, es digitalizar el fideicomiso, con la diferencia de que permite microinversiones sin necesidad de comprar una unidad entera. El mercado secundario, durante o al final del proyecto, depende de las cláusulas del \textit{white paper}.

\section{El desarrollista y la propiedad del inmueble}

El desarrollista debe demostrar, como condición, que es \textbf{propietario del inmueble}. No es necesariamente la constructora: puede contratar a un grupo de arquitectos o a un estudio de arquitectura. La constructora forma parte de la plataforma de tokenización. La información sobre la propiedad del inmueble es la que se traslada al token.

Las obligaciones del desarrollista (respetar las normas de construcción, presentar los planos) constarán en el \textit{white paper}, que es como la transcripción de un fideicomiso. Por ello se necesita conocer el contrato de fideicomiso inmobiliario, la clave asimétrica, la función de hash y el funcionamiento de \textit{blockchain} para entender la tokenización: se requiere un \textbf{doble conocimiento}. Esto complejiza el escenario jurídico para abogados, escribanos y jueces, cuando entre en alguna causal de disputa.

\section{Cuadro Cornell del capítulo}

\begin{longtable}{
    >{\raggedright\arraybackslash}p{0.27\textwidth}
    >{\raggedright\arraybackslash}p{0.65\textwidth}
}
\toprule
\textbf{Preguntas / palabras clave} & \textbf{Notas / respuestas} \\
\midrule
\endfirsthead
\toprule
\textbf{Preguntas / palabras clave} & \textbf{Notas / respuestas} \\
\midrule
\endhead
\bottomrule
\endfoot

\textbf{¿Qué son los nodos, el hash y el sello de tiempo?} &
Los nodos son copias del libro contable en tiempo real que generan consenso. El hash (64 caracteres) valida cada bloque y el sello de tiempo hace la transacción inmutable: nadie puede quitar el token a su titular. \\[0.5em]

\textbf{¿Cómo garantiza blockchain que el token es mío?} &
La adquisición queda registrada en la red, con hash y sello de tiempo, y no puede borrarse ni revertirse; por eso el desarrollador no podría desconocer la compra. \\[0.5em]

\textbf{¿Qué diferencia hay con las criptomonedas si se pierde el usuario y la clave?} &
En las criptomonedas se pierde todo; en la tokenización existe autenticación previa del usuario (\textit{onboarding}), por lo que usuario y clave pueden restablecerse. \\[0.5em]

\textbf{¿Qué es el token de calidad?} &
Modelo en el que los tokens se emiten a medida que hay adherentes; los no emitidos siguen siendo del desarrollista. El token contiene contrato, adhesión y controles de identidad. \\[0.5em]

\textbf{¿Puede el desarrollista licuar el valor del token?} &
No: no le conviene, el token es uno solo, ya está dividido y registrado en \textit{blockchain}, y nadie puede subdividirlo. El fraccionamiento se define en el \textit{white paper}. \\[0.5em]

\textbf{¿Qué debe demostrar el desarrollista?} &
Que es propietario del inmueble; esa información se traslada al token. \\

\midrule
\multicolumn{2}{p{0.92\textwidth}}{
\textbf{Resumen:} La red de \textit{blockchain} (nodos, consenso, hash y sello de tiempo) otorga inmutabilidad y fecha cierta a cada adquisición de token. En la tokenización argentina descrita, los tokens se emiten a medida que hay adherentes y el desarrollista conserva los no emitidos; no puede licuar ni subdividir tokens ya registrados. El \textit{onboarding}, que puede realizar el escribano, permite recuperar usuario y clave, a diferencia de las criptomonedas.
} \\
\end{longtable}

%==================================================================
% CAPÍTULO 8
%==================================================================
\chapter[Rol notarial, proyectos y hoja de ruta]{Rol del profesional: proyectos de tokenización, hoja de ruta y cierre}

\section{Ejemplos de proyectos}

Se mencionan los siguientes proyectos como ejemplos:

\begin{itemize}
    \item \textbf{Land Token}: fue uno de los primeros.
    \item \textbf{Agrotoken}.
    \item \textbf{Metro Futuro}: no trabaja en Argentina sino en España, pero es un desarrollo argentino.
    \item \textbf{Open Vino}: no es sobre inmuebles sino sobre botellas de vino; es una bodega.
    \item \textbf{Pala Blockchain}: empresa de tokenización, caso de éxito y con convenio firmado con el Colegio de Escribanos, que se vale del escribano para el proyecto de tokenización.
\end{itemize}

\section{Hoja de ruta (roadmap) del escribano en un proyecto de tokenización}

La hoja de ruta ordena los pasos que se seguirían al trabajar en un proyecto de tokenización.

\begin{enumerate}
    \item \textbf{Análisis documental} de todo el proyecto.
    \item \textbf{Propuesta de honorarios}; suscripción de un convenio de confidencialidad y un contrato.
    \item \textbf{Generación del \textit{white paper}} con el proyecto de fideicomiso; entrega a la empresa para su devolución y sugerencias.
    \item \textbf{Pedido de certificados} para la transmisión del inmueble al fideicomiso que procederá a la tokenización.
    \item \textbf{\textit{Closing}}: firma de todos los instrumentos públicos y privados.
    \item \textbf{Inscripción del fideicomiso} del inmueble: no se puede obviar, porque se tokeniza sobre un inmueble sobre el cual efectivamente se tiene el derecho de propiedad, aunque sea propiedad fiduciaria.
    \item \textbf{Tokenización}: términos y condiciones con los consumidores, comercialización mediante \textit{marketing} y plataforma inmobiliaria que se generará.
\end{enumerate}

\section{¿Sólo inmuebles? Las obligaciones negociables fueron lo primero}

La tokenización en Argentina no se limita a inmuebles. El proyecto de tokenización arranca con las \textbf{obligaciones negociables} que pueden emitir las empresas: es un negocio financiero. En Argentina, lo primero que se permitió tokenizar son las obligaciones negociables.

\section{Qué debe llevarse el estudiante}

No se pretende que el estudiante responda de memoria las resoluciones de la CNV, sino que:

\begin{itemize}
    \item tenga la \textbf{noción de tokenización};
    \item sepa que existe un \textbf{sandbox} y una \textbf{regulación} (de la CNV); y
    \item sepa \textbf{dónde buscar} cuando un cliente consulte, una vez finalizada la materia.
\end{itemize}

\begin{important}
No se exige memorizar las resoluciones. Lo que importa es conocer la noción de tokenización, saber que existe un \textit{sandbox} normativo y una regulación de la CNV, y saber dónde buscar cuando un cliente consulte.
\end{important}

\section{Cuadro Cornell del capítulo}

\begin{longtable}{
    >{\raggedright\arraybackslash}p{0.27\textwidth}
    >{\raggedright\arraybackslash}p{0.65\textwidth}
}
\toprule
\textbf{Preguntas / palabras clave} & \textbf{Notas / respuestas} \\
\midrule
\endfirsthead
\toprule
\textbf{Preguntas / palabras clave} & \textbf{Notas / respuestas} \\
\midrule
\endhead
\bottomrule
\endfoot

\textbf{¿Cuál es el rol del escribano en la tokenización?} &
Verificar la titularidad y gravámenes, realizar el \textit{onboarding} e identidad de los usuarios, certificar la adhesión al fideicomiso (por ejemplo, vía firma a distancia SDF 2.0) y acompañar todo el \textit{roadmap}: documentación, \textit{white paper}, \textit{closing}, inscripción y términos con consumidores. \\[0.5em]

\textbf{¿Qué pasos sigue el profesional en un proyecto?} &
Análisis documental; honorarios, confidencialidad y contrato; \textit{white paper} y devolución de la empresa; certificados para transmitir el inmueble al fideicomiso; \textit{closing}; inscripción del fideicomiso; tokenización (términos y condiciones, \textit{marketing} y plataforma). \\[0.5em]

\textbf{¿Qué se tokenizó primero en Argentina?} &
Las obligaciones negociables que pueden emitir las empresas. \\[0.5em]

\textbf{¿Qué se espera que se lleve el estudiante?} &
La noción de tokenización, saber que existe un \textit{sandbox} normativo y una regulación de la CNV, y saber dónde buscar cuando un cliente consulte; no se exige memorizar las resoluciones. \\

\midrule
\multicolumn{2}{p{0.92\textwidth}}{
\textbf{Resumen:} El profesional, y en especial el escribano, sigue siendo indispensable en la tokenización: verifica titularidad y gravámenes, identifica a las partes, redacta y lee el \textit{white paper}, escritura e inscribe el fideicomiso. La hoja de ruta ordena los pasos del proyecto, y lo esencial es comprender la noción de tokenización y saber dónde buscar la regulación.
} \\
\end{longtable}

%==================================================================
% SÍNTESIS FINAL
%==================================================================
\chapter[Síntesis final]{Síntesis final de la clase}

La clase unió dos temas que parecen lejanos pero comparten la misma lógica jurídica: la \textbf{donación con pacto de reversión} y la \textbf{tokenización inmobiliaria}. En ambos casos la base es el derecho real argentino (título, modo y registración, con escritura pública y efectos publicitarios del registro) y el límite lo pone el derecho constitucional de propiedad (art.~17 CN).

En la donación, el pacto de reversión protege al donante frente al orden de los fallecimientos, pero su duración más allá de los 10 años es discutible y, en la práctica, sólo se resuelve con sentencia o renuncia voluntaria. Por eso la planificación sucesoria exige mirar título por título (antes de 2015, entre 2015 y 2020 y después de 2020) y entender los vínculos familiares.

En la tokenización, el token no transmite el dominio sino un derecho personal de propiedad sobre un bien intangible vinculado a un activo real (RWA). Su funcionamiento se apoya en el fideicomiso, el \textit{white paper}, los \textit{smart contracts}, los oráculos y la \textit{blockchain} (nodos, hash y sello de tiempo), y su control se canaliza a través del \textit{sandbox} normativo de la CNV y de los PSAV inscriptos. Lo que cambia es la causa fuente del dominio, no las reglas del Código Civil y Comercial.

El profesional, y en especial el escribano, sigue siendo indispensable para verificar titularidad y gravámenes, identificar a las partes, redactar y leer el \textit{white paper}, escriturar y proteger al cliente de esquemas fraudulentos. La consigna general de la cátedra es no memorizar, sino \textbf{aplicar el Código con fundamento y saber defender una postura}.

\end{document}
